% ============================================================
% 高考数学满分突破 · 解析几何的代数化
% 编译命令：xelatex -interaction=nonstopmode 文件名.tex
% ============================================================
\documentclass[12pt, a4paper]{ctexart}
\usepackage[margin=2.2cm]{geometry}
\usepackage{amsmath, amssymb}
\usepackage{enumitem}
\usepackage{tcolorbox}
\usepackage{booktabs}
\usepackage{array}
\usepackage{fancyhdr}
\usepackage{xcolor}
\usepackage{tikz}
\usepackage{hyperref}
\usepackage{graphicx}
\usepackage{multirow}
\usepackage{pgfplots}
\pgfplotsset{compat=1.18}

% 颜色定义
\definecolor{mainblue}{RGB}{41, 98, 168}
\definecolor{lightblue}{RGB}{230, 240, 255}
\definecolor{tipgreen}{RGB}{34, 139, 87}
\definecolor{lightgreen}{RGB}{235, 250, 240}
\definecolor{warnred}{RGB}{200, 50, 50}
\definecolor{lightred}{RGB}{255, 240, 240}
\definecolor{examorange}{RGB}{210, 120, 20}
\definecolor{lightorange}{RGB}{255, 248, 235}
\definecolor{apppurple}{RGB}{120, 60, 170}
\definecolor{lightpurple}{RGB}{245, 238, 255}

% 页眉页脚
\pagestyle{fancy}
\fancyhf{}
\fancyhead[L]{\small\textcolor{mainblue}{\textbf{高考满分突破 · 解析几何的代数化}}}
\fancyhead[R]{\small\textcolor{mainblue}{高三冲刺讲义}}
\fancyfoot[C]{\small —\ \thepage\ —}
\renewcommand{\headrulewidth}{0.6pt}
\renewcommand{\headrule}{\hbox to\headwidth{\color{mainblue}\leaders\hrule height \headrulewidth\hfill}}
\renewcommand{\footrulewidth}{0pt}

% tcolorbox 样式
\newtcolorbox{knowledgebox}[1][]{%
  colback=lightblue, colframe=mainblue,
  fonttitle=\bfseries\large, title=#1,
  arc=3mm, boxrule=0.8pt,
  left=8pt, right=8pt, top=6pt, bottom=6pt
}

\newtcolorbox{tipbox}[1][]{%
  colback=lightgreen, colframe=tipgreen,
  fonttitle=\bfseries, title=#1,
  arc=2mm, boxrule=0.6pt,
  left=6pt, right=6pt, top=4pt, bottom=4pt
}

\newtcolorbox{warnbox}[1][]{%
  colback=lightred, colframe=warnred,
  fonttitle=\bfseries, title=#1,
  arc=2mm, boxrule=0.6pt,
  left=6pt, right=6pt, top=4pt, bottom=4pt
}

\newtcolorbox{examplebox}[1][]{%
  colback=lightorange, colframe=examorange,
  fonttitle=\bfseries, title=#1,
  arc=3mm, boxrule=0.8pt,
  left=8pt, right=8pt, top=6pt, bottom=6pt
}

\newtcolorbox{appbox}[1][]{%
  colback=lightpurple, colframe=apppurple,
  fonttitle=\bfseries, title=#1,
  arc=3mm, boxrule=0.8pt,
  left=8pt, right=8pt, top=6pt, bottom=6pt
}

\setCJKmainfont{Hiragino Sans GB}
\setCJKsansfont{Heiti SC}

\begin{document}

% ============ 标题页 ============
\begin{center}
{\Huge\bfseries\textcolor{mainblue}{高考数学满分突破}}\\[6pt]
{\LARGE\bfseries\textcolor{mainblue}{解析几何的代数化}}\\[12pt]
\rule{0.6\textwidth}{1pt}\\[8pt]
{\normalsize 适用：高三冲刺\quad |\quad 课时：3课时}
\end{center}
\vspace{10pt}

% ============================================================
%  第一部分：知识点讲解
% ============================================================
\section*{\textcolor{mainblue}{一、知识点讲解}}

% ---------- 1.1 曲线 = 方程的几何图像 ----------
\subsection*{1.1\quad 曲线 = 方程的几何图像}

\begin{knowledgebox}[核心原理]
解析几何的灵魂是\textbf{以代数方程描述几何曲线}：
\begin{itemize}[leftmargin=2em]
  \item 直线：$Ax+By+C=0$（一般式）或 $y=kx+b$（斜截式）。
  \item 圆：$(x-a)^2+(y-b)^2=r^2$，圆心 $(a,b)$，半径 $r$。
  \item 椭圆：$\dfrac{x^2}{a^2}+\dfrac{y^2}{b^2}=1$（$a>b>0$），焦点在 $x$ 轴。
  \item 双曲线：$\dfrac{x^2}{a^2}-\dfrac{y^2}{b^2}=1$，渐近线 $y=\pm\dfrac{b}{a}x$。
  \item 抛物线：$y^2=2px$（$p>0$），焦点 $\left(\dfrac{p}{2},0\right)$，准线 $x=-\dfrac{p}{2}$。
\end{itemize}
\end{knowledgebox}

\begin{knowledgebox}[方程与曲线的对应关系]
\begin{enumerate}[leftmargin=1.5em]
  \item 曲线上的点满足方程（\textbf{纯粹性}）。
  \item 满足方程的点都在曲线上（\textbf{完备性}）。
  \item 求曲线的交点 = 解联立方程组。
  \item 判断点与曲线的位置关系：代入方程，看符号。
\end{enumerate}
\end{knowledgebox}

\begin{warnbox}[易错提醒]
\begin{enumerate}[leftmargin=1.5em]
  \item 椭圆标准方程中 $a$ 是\textbf{半长轴}，$c^2=a^2-b^2$。双曲线中 $c^2=a^2+b^2$。
  \item 抛物线焦点弦长公式 $|AB|=x_1+x_2+p$ 只适用于焦点在 $x$ 轴的 $y^2=2px$ 型。
  \item "曲线过某点"意味着该点坐标代入方程成立。
\end{enumerate}
\end{warnbox}

% ---------- 1.2 齐次化方法的代数本质 ----------
\subsection*{1.2\quad 齐次化方法的代数本质}

\begin{knowledgebox}[什么是齐次化]
齐次化是将\textbf{非齐次方程}转化为\textbf{齐次方程}的技术，核心思想是\textbf{以一个变量作分母，消去常数项}。

设直线 $l:y=kx+m$ 与曲线 $f(x,y)=0$ 交于 $A,B$ 两点。
将 $y=kx+m$ 代入曲线方程，得到关于 $x$ 的二次方程。
\end{knowledgebox}

\begin{knowledgebox}[齐次化的标准操作]
\textbf{齐次化处理}：将曲线方程 $Ax^2+Bxy+Cy^2+Dx+Ey+F=0$ 与直线 $y=kx+m$ 联立，利用 $m=y-kx$ 替换，可以将方程化为关于 $\dfrac{x}{y}$ 或 $\dfrac{y}{x}$ 的齐次方程。

\medskip
\noindent\textbf{关键步骤}：
\begin{enumerate}[leftmargin=1.5em]
  \item 设 $A(x_1,y_1),B(x_2,y_2)$ 是直线与曲线的两个交点。
  \item 将 $m=y-kx$ 代入曲线方程，整理为 $Px^2+Qxy+Ry^2=0$ 的齐次形式。
  \item 两边除以 $x^2$（或 $y^2$），得到关于 $t=\dfrac{y}{x}$ 的二次方程 $Rt^2+Qt+P=0$。
  \item $t_1=\dfrac{y_1}{x_1},t_2=\dfrac{y_2}{x_2}$ 就是该方程的两个根。
\end{enumerate}
\end{knowledgebox}

\begin{tipbox}[齐次化的优势]
齐次化后可以直接用韦达定理得到 $\dfrac{y_1}{x_1}+\dfrac{y_2}{x_2}$ 和 $\dfrac{y_1}{x_1}\cdot\dfrac{y_2}{x_2}$，这在处理\textbf{斜率之和、斜率之积}类问题时极其高效，避免了繁琐的坐标运算。
\end{tipbox}

% ---------- 1.3 极点极线与射影背景 ----------
\subsection*{1.3\quad 极点极线与射影几何背景}

\begin{knowledgebox}[极点极线的定义]
对于椭圆 $\dfrac{x^2}{a^2}+\dfrac{y^2}{b^2}=1$，点 $P(x_0,y_0)$ 关于椭圆的\textbf{极线}方程为：
\[
  \frac{x_0 x}{a^2}+\frac{y_0 y}{b^2}=1.
\]
类似地，对于抛物线 $y^2=2px$，点 $P(x_0,y_0)$ 的极线为 $y_0 y=p(x+x_0)$。
\end{knowledgebox}

\begin{knowledgebox}[极点极线的核心性质]
\begin{enumerate}[leftmargin=1.5em]
  \item 若点 $P$ 在椭圆外，从 $P$ 引两条切线，切点分别为 $A,B$，则 $AB$ 就是 $P$ 的极线。
  \item 若过 $P$ 的直线交椭圆于 $M,N$ 两点，交极线于 $Q$，则 $P,Q$ 调和分割 $M,N$，即
  \[
    \frac{PM}{PN}=-\frac{QM}{QN},\quad\text{或}\quad\frac{1}{PM}+\frac{1}{PN}=\frac{2}{PQ}.
  \]
  \item 极线方程的形式与曲线方程\textbf{完全对称}：将曲线方程中的 $x^2\to x_0 x,\;y^2\to y_0 y$。
\end{enumerate}
\end{knowledgebox}

\begin{warnbox}[射影几何视角]
极点极线是射影几何的基本概念。高考中虽然不要求射影几何理论，但极点极线的代数结论可以直接使用。理解其背景有助于快速记忆公式。
\end{warnbox}

% ---------- 1.4 变换思想简化计算 ----------
\subsection*{1.4\quad 变换思想简化计算}

\begin{knowledgebox}[常用变换策略]
\begin{enumerate}[leftmargin=1.5em]
  \item \textbf{平移变换}：将曲线中心（或顶点）移至原点，简化方程。
  \item \textbf{旋转变物线}：当椭圆的对称轴不与坐标轴平行时，用旋转变换。
  \item \textbf{参数方程}：用参数表示曲线上的点，将几何条件转化为参数关系。
  \item \textbf{设而不求}：设交点坐标但不具体解出，利用韦达定理整体处理。
  \item \textbf{齐次化}：将非齐次问题转化为齐次问题（详见1.2节）。
\end{enumerate}
\end{knowledgebox}

\begin{tipbox}[设而不求的核心公式]
设直线 $y=kx+m$ 与曲线交于 $A(x_1,y_1),B(x_2,y_2)$，联立后得二次方程 $ax^2+bx+c=0$，则
\[
  x_1+x_2=-\frac{b}{a},\quad x_1 x_2=\frac{c}{a},\quad |x_1-x_2|=\frac{\sqrt{b^2-4ac}}{|a|}.
\]
弦长公式：$|AB|=\sqrt{1+k^2}\,|x_1-x_2|=\sqrt{1+k^2}\cdot\dfrac{\sqrt{\Delta}}{|a|}$。
\end{tipbox}

% ============================================================
%  第二部分：例题精讲
% ============================================================
\section*{\textcolor{examorange}{二、例题精讲}}

\begin{examplebox}[【例题1】联立方程求交点]
求直线 $y=2x+1$ 与椭圆 $\dfrac{x^2}{4}+y^2=1$ 的交点坐标。
\end{examplebox}

\noindent\textbf{解}：将 $y=2x+1$ 代入椭圆方程：
\[
  \frac{x^2}{4}+(2x+1)^2=1\;\Longrightarrow\;\frac{x^2}{4}+4x^2+4x+1=1\;\Longrightarrow\;\frac{17x^2}{4}+4x=0.
\]
\[
  x\!\left(\frac{17x}{4}+4\right)=0\;\Longrightarrow\;x_1=0,\;x_2=-\frac{16}{17}.
\]
对应 $y_1=1,\;y_2=2\times\!\left(-\dfrac{16}{17}\right)+1=-\dfrac{15}{17}$。

交点为 $(0,1)$ 和 $\left(-\dfrac{16}{17},-\dfrac{15}{17}\right)$。

\vspace{6pt}

\begin{examplebox}[【例题2】韦达定理与弦长]
已知直线 $l:y=x+m$ 与椭圆 $\dfrac{x^2}{4}+\dfrac{y^2}{3}=1$ 交于 $A,B$ 两点，若 $|AB|=3$，求 $m$ 的值。
\end{examplebox}

\noindent\textbf{解}：联立 $y=x+m$ 与 $\dfrac{x^2}{4}+\dfrac{y^2}{3}=1$：
\[
  3x^2+4(x+m)^2=12\;\Longrightarrow\;7x^2+8mx+4m^2-12=0.
\]
设 $A(x_1,y_1),B(x_2,y_2)$，则 $x_1+x_2=-\dfrac{8m}{7},\;x_1x_2=\dfrac{4m^2-12}{7}$。

弦长：
\[
  |AB|=\sqrt{1+1^2}\,|x_1-x_2|=\sqrt{2}\cdot\frac{\sqrt{(x_1+x_2)^2-4x_1x_2}}{1}=\sqrt{2}\cdot\frac{\sqrt{\dfrac{64m^2}{49}-\dfrac{16m^2-48}{7}}}{1}.
\]
化简：$|AB|^2=2\cdot\dfrac{64m^2-49(16m^2-48)/7}{49}=2\cdot\dfrac{64m^2-112m^2+336}{49}=2\cdot\dfrac{336-48m^2}{49}$。

由 $|AB|=3$：$9=\dfrac{2(336-48m^2)}{49}$，解得 $48m^2=336-\dfrac{441}{2}=\dfrac{231}{2}$，即 $m^2=\dfrac{231}{96}=\dfrac{77}{32}$，$m=\pm\dfrac{\sqrt{77}}{4\sqrt{2}}=\pm\dfrac{\sqrt{154}}{8}$。

\vspace{6pt}

\begin{examplebox}[【例题3】齐次化处理斜率之和]
过点 $P(2,0)$ 的直线 $l$ 与椭圆 $\dfrac{x^2}{4}+y^2=1$ 交于 $A,B$ 两点，求直线 $PA,PB$ 斜率之和 $k_{PA}+k_{PB}$。
\end{examplebox}

\noindent\textbf{解}：设 $l:x=ty+2$（避免斜率不存在的情况），代入椭圆方程：
\[
  \frac{(ty+2)^2}{4}+y^2=1\;\Longrightarrow\;t^2y^2+4ty+4+4y^2=4\;\Longrightarrow\;(t^2+4)y^2+4ty=0.
\]
设 $A(x_1,y_1),B(x_2,y_2)$，则 $y_1+y_2=-\dfrac{4t}{t^2+4},\;y_1y_2=0$。

由于 $y_1y_2=0$，不妨设 $y_1=0$，则 $A(2,0)=P$。这说明 $l$ 过 $P$ 时 $A=P$，需要调整策略。

\noindent\textbf{改用齐次化}：将椭圆方程 $\dfrac{x^2}{4}+y^2=1$ 中的 $1$ 用 $\dfrac{x-2}{-2}\cdot(-1)+\dfrac{y}{1}\cdot0$ 替换不合适。更直接的方法是设 $A(x_1,y_1),B(x_2,y_2)$，$l:y=k(x-2)$，代入椭圆：
\[
  (1+4k^2)x^2-16k^2x+16k^2-4=0.
\]
$x_1+x_2=\dfrac{16k^2}{1+4k^2},\;x_1x_2=\dfrac{16k^2-4}{1+4k^2}$。

\[
  k_{PA}+k_{PB}=\frac{y_1}{x_1-2}+\frac{y_2}{x_2-2}=\frac{k(x_1-2)}{x_1-2}+\frac{k(x_2-2)}{x_2-2}=k+k=2k.
\]
但注意 $P$ 是直线上的定点，$A,B$ 均在过 $P$ 的直线上，故 $k_{PA}=k_{PB}=k$，其和为 $2k$。

\begin{tipbox}[齐次化的正确使用场景]
齐次化最适合处理"\textbf{直线过曲线外某点}，求斜率之和/积"的问题。当直线过曲线\textbf{上}的点时，直接计算更简便。
\end{tipbox}

\vspace{6pt}

\begin{examplebox}[【例题4】齐次化——标准题型]
已知椭圆 $\dfrac{x^2}{2}+y^2=1$，$A(-2,0)$。过 $A$ 的直线 $l$ 交椭圆于 $M,N$ 两点（$A$ 不在椭圆上），求 $k_{AM}\cdot k_{AN}$ 的取值范围。
\end{examplebox}

\noindent\textbf{解}：设 $l:x=my-2$，代入椭圆：
\[
  \frac{(my-2)^2}{2}+y^2=1\;\Longrightarrow\;(m^2+2)y^2-4my+2=0.
\]
设 $M(x_1,y_1),N(x_2,y_2)$，则 $y_1+y_2=\dfrac{4m}{m^2+2},\;y_1y_2=\dfrac{2}{m^2+2}$。

\[
  k_{AM}\cdot k_{AN}=\frac{y_1}{x_1+2}\cdot\frac{y_2}{x_2+2}=\frac{y_1y_2}{(x_1+2)(x_2+2)}.
\]
由 $x_i=my_i-2$，得 $x_i+2=my_i$，故
\[
  k_{AM}\cdot k_{AN}=\frac{y_1y_2}{m^2 y_1 y_2}=\frac{1}{m^2}.
\]
$m$ 取任意实数，$m^2>0$，故 $k_{AM}\cdot k_{AN}=\dfrac{1}{m^2}>0$。

但需确保 $\Delta>0$：$16m^2-8(m^2+2)>0\Rightarrow m^2>2$，故 $\dfrac{1}{m^2}\in\left(0,\dfrac{1}{2}\right)$。

\vspace{6pt}

\begin{examplebox}[【例题5】极点极线——代数应用]
已知椭圆 $\dfrac{x^2}{4}+y^2=1$，$P(1,1)$。设过 $P$ 的直线交椭圆于 $A,B$ 两点，求 $PA$ 中点 $M$ 的轨迹方程。
\end{examplebox}

\noindent\textbf{解}：设 $A(x_1,y_1),B(x_2,y_2)$，中点 $M(x_0,y_0)$。则 $x_1+x_2=2x_0,\;y_1+y_2=2y_0$。

$A,B$ 在椭圆上：$\dfrac{x_i^2}{4}+y_i^2=1$，两式相减：
\[
  \frac{x_1^2-x_2^2}{4}+(y_1^2-y_2^2)=0\;\Longrightarrow\;\frac{(x_1+x_2)(x_1-x_2)}{4}+(y_1+y_2)(y_1-y_2)=0.
\]
\[
  \frac{2x_0}{4}(x_1-x_2)+2y_0(y_1-y_2)=0\;\Longrightarrow\;k_{AB}=\frac{y_1-y_2}{x_1-x_2}=-\frac{x_0}{4y_0}.
\]
又 $A,B,P$ 共线，$k_{AP}=\dfrac{y_0-1}{x_0-1}=k_{AB}=-\dfrac{x_0}{4y_0}$，故
\[
  (y_0-1)\cdot4y_0=-x_0(x_0-1)\;\Longrightarrow\;4y_0^2-4y_0=-x_0^2+x_0.
\]
\[
  x_0^2+4y_0^2-x_0-4y_0=0\;\Longrightarrow\;\left(x_0-\frac{1}{2}\right)^2+4\!\left(y_0-\frac{1}{2}\right)^2=\frac{5}{4}.
\]
$PA$ 中点轨迹为椭圆 $\dfrac{\left(x-\frac{1}{2}\right)^2}{\frac{5}{4}}+\dfrac{\left(y-\frac{1}{2}\right)^2}{\frac{5}{16}}=1$ 在原椭圆内部的部分。

\vspace{6pt}

\begin{examplebox}[【例题6】极线方程的直接应用]
已知椭圆 $\dfrac{x^2}{9}+\dfrac{y^2}{4}=1$，点 $P(6,2)$。求 $P$ 关于椭圆的极线方程。
\end{examplebox}

\noindent\textbf{解}：直接代入极线公式 $\dfrac{x_0 x}{a^2}+\dfrac{y_0 y}{b^2}=1$：
\[
  \frac{6x}{9}+\frac{2y}{4}=1\;\Longrightarrow\;\frac{2x}{3}+\frac{y}{2}=1\;\Longrightarrow\;4x+3y=6.
\]
极线方程为 $4x+3y=6$。

\begin{tipbox}[规律总结]
极线公式可以直接由曲线方程"替换"得到：$x^2\to x_0 x,\;y^2\to y_0 y$。记忆口诀：\textbf{曲线方程中平方项换成点坐标乘以动点坐标}。
\end{tipbox}

\vspace{6pt}

\begin{examplebox}[【例题7】参数方程简化计算]
椭圆 $\dfrac{x^2}{4}+y^2=1$ 上一动点 $P$，求 $P$ 到直线 $x+y-4=0$ 的最短距离。
\end{examplebox}

\noindent\textbf{解}：设 $P(2\cos\theta,\sin\theta)$，则
\[
  d=\frac{|2\cos\theta+\sin\theta-4|}{\sqrt{2}}.
\]
$2\cos\theta+\sin\theta=\sqrt{5}\sin(\theta+\varphi)$（其中 $\tan\varphi=2$），取值范围 $[-\sqrt{5},\sqrt{5}]$。

故 $|2\cos\theta+\sin\theta-4|$ 的最小值为 $4-\sqrt{5}$（因为 $\sqrt{5}<4$），最短距离为
\[
  d_{\min}=\frac{4-\sqrt{5}}{\sqrt{2}}=\frac{(4-\sqrt{5})\sqrt{2}}{2}.
\]

\begin{tipbox}[方法总结]
椭圆上点的参数方程 $(a\cos\theta,b\sin\theta)$ 将二元问题转化为一元三角函数问题，是求最值的常用技巧。
\end{tipbox}

\vspace{6pt}

\begin{examplebox}[【例题8】设而不求——整体代入]
直线 $y=kx+1$ 与双曲线 $x^2-y^2=1$ 交于 $A,B$ 两点，若 $\overrightarrow{OA}\cdot\overrightarrow{OB}=0$（$O$ 为原点），求 $k$ 的值。
\end{examplebox}

\noindent\textbf{解}：联立 $y=kx+1$ 与 $x^2-y^2=1$：
\[
  x^2-(kx+1)^2=1\;\Longrightarrow\;(1-k^2)x^2-2kx-2=0.
\]
设 $A(x_1,y_1),B(x_2,y_2)$，则 $x_1+x_2=\dfrac{2k}{1-k^2},\;x_1x_2=\dfrac{-2}{1-k^2}$（需 $k^2\neq1$ 且 $\Delta>0$）。

$\overrightarrow{OA}\cdot\overrightarrow{OB}=x_1x_2+y_1y_2=0$，而 $y_i=kx_i+1$：
\[
  y_1y_2=(kx_1+1)(kx_2+1)=k^2x_1x_2+k(x_1+x_2)+1.
\]
\[
  x_1x_2+y_1y_2=(1+k^2)x_1x_2+k(x_1+x_2)+1=(1+k^2)\cdot\frac{-2}{1-k^2}+k\cdot\frac{2k}{1-k^2}+1.
\]
\[
  =\frac{-2-2k^2+2k^2}{1-k^2}+1=\frac{-2}{1-k^2}+1=\frac{k^2-3}{1-k^2}=0.
\]
故 $k^2=3$，$k=\pm\sqrt{3}$。验证 $\Delta>0$ 且 $k^2\neq1$，成立。

\vspace{6pt}

\begin{examplebox}[【例题9】抛物线焦点弦问题]
抛物线 $y^2=4x$，过焦点 $F(1,0)$ 的直线交抛物线于 $A,B$ 两点，若 $|AF|=3$，求 $|BF|$。
\end{examplebox}

\noindent\textbf{解}：设 $A(x_1,y_1),B(x_2,y_2)$。抛物线 $y^2=4x$ 的 $p=2$。

焦点弦公式：$\dfrac{1}{|AF|}+\dfrac{1}{|BF|}=\dfrac{2}{p}=1$。

由 $|AF|=3$：$\dfrac{1}{3}+\dfrac{1}{|BF|}=1$，解得 $|BF|=\dfrac{3}{2}$。

\begin{tipbox}[规律总结]
抛物线焦点弦的倒数和公式：$\dfrac{1}{|AF|}+\dfrac{1}{|BF|}=\dfrac{2}{p}$，这是一个非常简洁的结论，记忆方便且计算快速。
\end{tipbox}

\vspace{6pt}

\begin{examplebox}[【例题10】平移变换简化方程]
将椭圆 $\dfrac{x^2}{4}+y^2=1$ 向右平移 $1$ 个单位，再向上平移 $2$ 个单位，求平移后的方程。
\end{examplebox}

\noindent\textbf{解}：平移变换：$x\to x-1,\;y\to y-2$。
\[
  \frac{(x-1)^2}{4}+(y-2)^2=1.
\]
\begin{tipbox}[变换思想]
平移变换的本质是坐标替换：将原方程中的 $x$ 替换为 $x-h$，$y$ 替换为 $y-k$（$(h,k)$ 为平移向量）。这个操作在高考中经常用于将非标准位置的曲线转化为标准形式。
\end{tipbox}

\vspace{6pt}

\begin{examplebox}[【例题11】齐次化——斜率之积为定值]
椭圆 $\dfrac{x^2}{4}+\dfrac{y^2}{3}=1$，过点 $T(4,0)$ 的直线 $l$ 交椭圆于 $A,B$。设 $k_1,k_2$ 分别为直线 $CA,CB$ 的斜率（$C$ 为椭圆右顶点 $(2,0)$），证明 $k_1\cdot k_2$ 为定值。
\end{examplebox}

\noindent\textbf{证}：设 $l:x=my+4$，代入椭圆：
\[
  3(my+4)^2+4y^2=12\;\Longrightarrow\;(3m^2+4)y^2+24my+36=0.
\]
设 $A(x_1,y_1),B(x_2,y_2)$，则 $y_1+y_2=-\dfrac{24m}{3m^2+4},\;y_1y_2=\dfrac{36}{3m^2+4}$。

\[
  k_1\cdot k_2=\frac{y_1}{x_1-2}\cdot\frac{y_2}{x_2-2}=\frac{y_1y_2}{(x_1-2)(x_2-2)}.
\]
$x_i=my_i+4$，故 $x_i-2=my_i+2$，$(x_1-2)(x_2-2)=m^2y_1y_2+2m(y_1+y_2)+4$。

代入：
\[
  (x_1-2)(x_2-2)=m^2\cdot\frac{36}{3m^2+4}+2m\cdot\frac{-24m}{3m^2+4}+4=\frac{36m^2-48m^2+12m^2+16}{3m^2+4}=\frac{16}{3m^2+4}.
\]
\[
  k_1\cdot k_2=\frac{\frac{36}{3m^2+4}}{\frac{16}{3m^2+4}}=\frac{36}{16}=\frac{9}{4}.
\]
$k_1\cdot k_2=\dfrac{9}{4}$ 为定值，证毕。

\vspace{6pt}

\begin{examplebox}[【例题12】综合——面积与弦长]
椭圆 $\dfrac{x^2}{4}+y^2=1$，直线 $l:y=kx+m$ 与椭圆交于 $A,B$ 两点，$\triangle AOB$（$O$ 为原点）面积最大时，求 $k$ 与 $m$ 的关系。
\end{examplebox}

\noindent\textbf{解}：联立得 $(1+4k^2)x^2+8kmx+4m^2-4=0$。

$S_{\triangle AOB}=\dfrac{1}{2}|AB|\cdot d(O,l)=\dfrac{1}{2}\cdot\dfrac{\sqrt{1+k^2}\cdot\sqrt{\Delta}}{1+4k^2}\cdot\dfrac{|m|}{\sqrt{1+k^2}}=\dfrac{|m|\sqrt{\Delta}}{2(1+4k^2)}$。

其中 $\Delta=64k^2m^2-4(1+4k^2)(4m^2-4)=16(1+4k^2-m^2)$。

\[
  S=\frac{|m|\cdot4\sqrt{1+4k^2-m^2}}{2(1+4k^2)}=\frac{2|m|\sqrt{1+4k^2-m^2}}{1+4k^2}.
\]
令 $u=m^2$，$S^2=\dfrac{4u(1+4k^2-u)}{(1+4k^2)^2}$。对 $u$ 求导取极值：$u=\dfrac{1+4k^2}{2}$，即 $m^2=\dfrac{1+4k^2}{2}$ 时面积最大。

此时 $S_{\max}=\dfrac{2\cdot\frac{1+4k^2}{2}}{1+4k^2}=1$（需验证 $\Delta>0$，即 $m^2<1+4k^2$，满足）。

% ============================================================
%  第三部分：应用题
% ============================================================
\section*{\textcolor{apppurple}{三、应用题}}

\begin{appbox}[【应用1】卫星轨道]
某卫星的轨道是椭圆，地球（视为一个焦点）位于椭圆的一个焦点上。已知轨道的长半轴 $a=10000$ km，离心率 $e=0.2$。求卫星到地球的最近距离和最远距离。
\end{appbox}

\noindent\textbf{解}：$c=ae=2000$ km，短半轴 $b=\sqrt{a^2-c^2}=\sqrt{10^8-4\times10^6}=4000\sqrt{6}\approx9798$ km。

近地点距离 $=a-c=8000$ km，远地点距离 $=a+c=12000$ km。

\vspace{6pt}

\begin{appbox}[【应用2】抛物面天线]
一个抛物面天线的截面是抛物线，开口向上，顶点在原点。已知天线深度为 $0.5$ m，口径为 $2$ m。求焦点位置（即接收器应放置的位置）。
\end{appbox}

\noindent\textbf{解}：设抛物线方程 $x^2=2py$。口径 $2$ m 即 $x=\pm1$ 时 $y=0.5$。
\[
  1=2p\times0.5=p\;\Longrightarrow\;p=1.
\]
焦点在 $\left(0,\dfrac{p}{2}\right)=(0,0.5)$，即接收器应在天线顶点上方 $0.5$ m 处。

\vspace{6pt}

\begin{appbox}[【应用3】导航定位]
两个信号站 $F_1(-3,0),F_2(3,0)$（单位：km）同时发出信号。某船收到 $F_1$ 的信号比 $F_2$ 的信号晚了 $c$ 微秒（信号传播速度 $v=300$ m/$\mu$s）。若船到两站距离之差恰好为 $4$ km，求船所在的双曲线方程。
\end{appbox}

\noindent\textbf{解}：距离差 $2a=4$，故 $a=2$。$c=3$，$b^2=c^2-a^2=9-4=5$。

双曲线方程：$\dfrac{x^2}{4}-\dfrac{y^2}{5}=1$。

信号延迟时间：$c$ 微秒对应距离差 $300c$ m $=0.3c$ km。若 $0.3c=4$，则 $c=\dfrac{40}{3}\approx13.3$ $\mu$s。

% ============================================================
%  第四部分：课堂练习
% ============================================================
\section*{\textcolor{tipgreen}{四、课堂练习}}

\begin{enumerate}[leftmargin=2em]
  \item 求椭圆 $\dfrac{x^2}{9}+\dfrac{y^2}{4}=1$ 的焦点坐标和离心率。

  \item 直线 $y=x-1$ 与抛物线 $y^2=4x$ 交于 $A,B$ 两点，求 $|AB|$。

  \item 已知双曲线 $\dfrac{x^2}{4}-\dfrac{y^2}{3}=1$，求其渐近线方程。

  \item 过椭圆 $\dfrac{x^2}{5}+y^2=1$ 的左焦点 $F_1$ 的直线交椭圆于 $A,B$，求 $\triangle ABF_2$（$F_2$ 为右焦点）周长的最小值。

  \item 点 $P(4,3)$ 关于椭圆 $\dfrac{x^2}{16}+\dfrac{y^2}{9}=1$ 的极线方程。

  \item 设椭圆 $\dfrac{x^2}{2}+y^2=1$，直线 $l$ 过点 $M(0,1)$ 交椭圆于 $A,B$，求 $\overrightarrow{MA}\cdot\overrightarrow{MB}$ 的取值范围。

  \item 抛物线 $y^2=8x$，焦点 $F$，过 $F$ 的直线交抛物线于 $A,B$，若 $|AF|=4$，求 $|BF|$ 及 $\triangle AOB$ 面积。

  \item 已知椭圆 $\dfrac{x^2}{4}+\dfrac{y^2}{3}=1$，点 $A(1,1)$ 在椭圆内部。过 $A$ 的弦 $MN$ 以 $A$ 为中点，求直线 $MN$ 的方程。
\end{enumerate}

\vspace{12pt}

\begin{knowledgebox}[参考答案]
\begin{enumerate}[leftmargin=2em]
  \item 焦点 $(\pm\sqrt{5},0)$，离心率 $e=\dfrac{\sqrt{5}}{3}$。

  \item 联立得 $x^2-6x+1=0$，$x_1+x_2=6,\;x_1x_2=1$。$|AB|=\sqrt{2}\cdot\sqrt{36-4}=4\sqrt{2}$。

  \item 渐近线：$y=\pm\dfrac{\sqrt{3}}{2}x$。

  \item $\triangle ABF_2$ 周长 $=|AF_1|+|AF_2|+|BF_1|+|BF_2|=4a+|AF_1|+|BF_1|=4\sqrt{5}+|AB|$。当 $AB$ 为通径（垂直于长轴）时 $|AB|$ 最小，通径长 $=\dfrac{2b^2}{a}=\dfrac{2}{\sqrt{5}}$。周长最小值 $=4\sqrt{5}+\dfrac{2}{\sqrt{5}}=\dfrac{22\sqrt{5}}{5}$。

  \item 极线：$\dfrac{4x}{16}+\dfrac{3y}{9}=1\Rightarrow\dfrac{x}{4}+\dfrac{y}{3}=1\Rightarrow 3x+4y=12$。

  \item 设 $l:y=kx+1$，代入椭圆得 $(1+2k^2)x^2+4kx=0$。$x_1=0$（即 $A$ 为 $M$ 本身时退化），$x_2=\dfrac{-4k}{1+2k^2}$。$\overrightarrow{MA}=x_1\cdot x_2+y_1\cdot y_2$，其中 $\overrightarrow{MA}=(x_1,kx_1),\overrightarrow{MB}=(x_2,kx_2)$。$\overrightarrow{MA}\cdot\overrightarrow{MB}=x_1x_2(1+k^2)=0$（因 $x_1=0$）。对于 $M$ 不在椭圆上的情况，用韦达定理：$x_1x_2=\dfrac{-2}{1+2k^2}$（$M$ 代入椭圆不成立时），$\overrightarrow{MA}\cdot\overrightarrow{MB}=\dfrac{-2(1+k^2)}{1+2k^2}\in[-2,-1)$。

  \item $|AF|+|BF|=\dfrac{2p}{\sin^2\alpha}$（焦点弦公式），但更直接：$\dfrac{1}{|AF|}+\dfrac{1}{|BF|}=\dfrac{2}{p}=1$，$|AF|=4\Rightarrow|BF|=\dfrac{4}{3}$。$S_{\triangle AOB}=\dfrac{1}{2}\cdot p\cdot|y_1-y_2|$。设 $A(4,y_A)$，由 $|AF|=4$ 且 $p=4$：$x_A+2=4\Rightarrow x_A=2,\;y_A=\pm4$。设 $A(2,4)$，$B\!\left(\dfrac{2}{3},\dfrac{-4}{3}\right)$。$S=\dfrac{1}{2}\cdot4\cdot\left|4-\!\left(-\dfrac{4}{3}\right)\right|=2\cdot\dfrac{16}{3}=\dfrac{16}{3}\sqrt{3}$（需取 $|y_1-y_2|$ 并调整为精确值 $S=\dfrac{16\sqrt{3}}{3}$）。

  \item 点差法：$\dfrac{x_1^2-x_2^2}{4}+y_1^2-y_2^2=0$，$\dfrac{(x_1+x_2)(x_1-x_2)}{4}+(y_1+y_2)(y_1-y_2)=0$。$x_1+x_2=2,\;y_1+y_2=2$，$k=\dfrac{y_1-y_2}{x_1-x_2}=-\dfrac{1}{4}$。直线 $MN:y-1=-\dfrac{1}{4}(x-1)$，即 $x+4y=5$。
\end{enumerate}
\end{knowledgebox}

\vfill
\begin{center}
\small\textcolor{gray}{本讲义仅供教学参考，例题可根据学生水平调整难度。}
\end{center}

\end{document}
